\documentclass[12pt]{article}

\usepackage[margin=1in]{geometry}
\hyphenpenalty=10000
\exhyphenpenalty=10000
\usepackage{times}
\usepackage{setspace}
\usepackage{titlesec}
\usepackage{hyperref}
\usepackage{amsmath}
\usepackage{booktabs}
\usepackage{parskip}
\usepackage[super,sort&compress]{natbib}
\usepackage{xcolor}
\usepackage{array}

\hypersetup{
    colorlinks=true,
    linkcolor=black,
    citecolor=black,
    urlcolor=blue
}

\titleformat{\section}{\large\bfseries}{}{0em}{}[\titlerule]
\titleformat{\subsection}{\normalsize\bfseries}{}{0em}{}

\title{
    \vspace{-1cm}
    {\large Boston University, Department of Computing \& Data Sciences}\\[0.5em]
    {\large M.S. Thesis Proposal}\\[1em]
    \textbf{What Does a Physics Network Actually Learn?\\
    Mechanistic Interpretability Across a Spectrum\\
    of Physics-Informed Neural Architectures}\\[1em]
}
\author{
    Achuthan Rathinam\\
    \small CDS DS 799, Master's Thesis\\[0.5em]
    \small \textbf{Thesis Advisor:} Siddharth Mishra-Sharma\\
    \small Department of Computing \& Data Sciences / Department of Physics
}
\date{April 2026}

\begin{document}

\maketitle
\thispagestyle{empty}

\onehalfspacing

% -------------------------------------------------------
\section{Problem Statement}

Physics-informed neural networks encode physical knowledge in two fundamentally
different ways. Some architectures are \textit{symmetry-constrained}: they are
designed so that outputs obey known symmetries (e.g., rotational equivariance)
regardless of what the network has learned internally. Others are
\textit{physics-explicit}: they directly parameterize a physical quantity such
as an energy or Hamiltonian, enforcing conservation laws through the structure
of the computation itself. Both approaches achieve strong empirical performance
on physical dynamics tasks. But a foundational question has not been answered:

\begin{center}
\textit{Do physics-explicit and symmetry-constrained architectures develop the same internal representations when trained on the same physical system, or do they discover fundamentally different computational strategies to produce the right answer?}
\end{center}

More concretely: when a network trained on Coulomb dynamics produces accurate
predictions, has it learned Coulomb's law, or has it found a different pattern
that happens to generalize within the training distribution? Symmetry
constraints guarantee that outputs transform correctly under rotation.
Physics-explicit structure guarantees energy conservation. Neither guarantee
tells us whether the network's hidden representations correspond to physically
meaningful quantities (forces, charge interactions, distance-dependent
potentials).

This thesis applies mechanistic interpretability methods to directly examine,
compare, and characterize the internal representations of five neural
architectures spanning a spectrum from physics-explicit to fully unconstrained,
all trained on an identical Coulomb particle dynamics task where ground-truth
physical quantities are analytically computable at every point.

% -------------------------------------------------------
\section{Background}

\subsection{A Spectrum of Physics Priors}

Modern physics-informed neural networks occupy a spectrum based on how deeply
physical knowledge is encoded:

\begin{center} \renewcommand{\arraystretch}{1.3}
  \begin{tabular}{>{\bfseries}lll} \toprule Tier & Architecture & Physics Prior
    \\ \midrule Physics-explicit & Hamiltonian Neural Network & Learns energy
    $H(q,p)$ directly; \\ & & conservation enforced by construction \\ \midrule
    Symmetry-        & EGNN                       & Equivariant message passing
    \\ constrained      & Attention EGNN             & Equivariant + learned
    interaction weights \\ & Equivariant Transformer    & Equivariant attention
  (Equiformer) \\ \midrule Unconstrained    & Non-equivariant GNN        &
Permutation symmetry only \\ & MLP                        & No structural prior
\\ \bottomrule \end{tabular} \end{center}

\subsection{Equivariant Graph Neural Networks}

The EGNN \cite{satorras2021en} performs message passing with coordinate updates
    constructed from equivariant operations on relative position vectors,
    ensuring $f(Rx, Rv) = R \cdot f(x,v)$ for all $R \in \mathrm{SO}(3)$. Each
    layer computes: \begin{align} m_{ij} &= \phi_e(h_i,\, h_j,\, \|x_i -
    x_j\|^2,\, a_{ij}) \\ x_i'   &= x_i + \sum_{j \neq i} (x_i -
  x_j)\,\phi_x(m_{ij}) \\ h_i' &= \phi_h\!\left(h_i,\, \textstyle\sum_{j \neq i}
m_{ij}\right) \end{align} where $h_i$ are node features, $x_i$ positions, and
$\phi_e, \phi_x, \phi_h$ learned MLPs.

\subsection{Hamiltonian Neural Networks}

Hamiltonian Neural Networks (HNNs) \cite{greydanus2019hamiltonian} learn a
scalar function $\mathcal{H}_\theta(q, p)$ and define dynamics via Hamilton's
equations: \[ \dot{q} = \frac{\partial \mathcal{H}_\theta}{\partial p}, \qquad
\dot{p} = -\frac{\partial \mathcal{H}_\theta}{\partial q} \] The loss penalizes
violations of these equations, so the network is explicitly trained to model
conserved energy. For Coulomb dynamics, a well-trained HNN should recover
$\mathcal{H} = \sum_i \frac{p_i^2}{2m} + \sum_{i < j} \frac{q_i q_j}{\|x_i -
x_j\|}$.

\subsection{Equivariant Transformers}

Equiformer \cite{liao2022equiformer} replaces message passing with equivariant
attention over irreducible representations of SO(3), achieving the same symmetry
guarantees as EGNN through a different computational mechanism. This enables a
direct test of whether symmetry-constrained architectures converge to similar
internal representations regardless of their specific implementation.

\subsection{Mechanistic Interpretability}

Mechanistic interpretability \cite{olah2020zoom, nanda2023progress} aims to
reverse-engineer the algorithms neural networks implement. This thesis uses four
complementary tools: \begin{itemize} \item \textbf{Linear probing:} Train
  lightweight classifiers on internal activations to test whether specific
  physical quantities are linearly encoded at each layer. \item \textbf{Causal
  intervention / activation patching:} Corrupt a specific input (e.g., flip
  charge $q_i$) and trace which components causally carry that change to the
  output. \item \textbf{Probing + ablation loop:} After identifying a probed
  feature (e.g., distance encoded at layer 2), ablate that direction in
  representation space and measure the resulting performance drop, upgrading
  from ``the model encodes $X$'' to ``the model \textit{uses} $X$ to make
  predictions.'' \item \textbf{Representational similarity (CKA):} Compare the
  geometry of internal representations across architectures to test whether
different models develop structurally equivalent solutions. \end{itemize}

\subsection{The Coulomb Dynamics Testbed}

$N$ particles with charges $q_i \in \{-1, +1\}$ interact via Coulomb forces: \[
F_{ij} = q_i q_j \frac{x_i - x_j}{\|x_i - x_j\|^3} \] Given initial positions
and velocities, the task is to predict positions after $T = 1000$ integration
steps. The governing law is fully known, providing ground-truth physical
quantities (pairwise distances, charge products, force magnitudes and
directions, total energy) against which internal representations can be
directly evaluated. This is the key advantage of this testbed over molecular or
protein datasets: there is no ambiguity about what the model ``should'' be
computing.

% -------------------------------------------------------
\section{Proposed Research}

\subsection{Central Thesis Claim}

We hypothesize that physics-explicit inductive bias (HNN) and symmetry-based
inductive bias (EGNN, Equiformer) both improve generalization, but produce
\textit{fundamentally different} internal representations. Specifically: only
physics-explicit architectures will develop hidden representations that directly
correspond to physical law (i.e., probes for Coulomb potential and energy will
succeed at intermediate layers), while symmetry-constrained architectures will
encode physically interpretable geometry (distances, relative positions) without
necessarily encoding the force law itself.

\subsection{Research Questions}

\begin{enumerate} \item Do internal representations encode physically
  interpretable quantities (pairwise distance, charge-product sign, force
  magnitude, force direction, potential energy), and at which layers do these
  emerge across architectures? \item Are probed representations causally active
  (does ablating a probed direction degrade performance, or is it a correlate
  rather than a cause?) \item Do equivariant architectures with different
  mechanisms (EGNN vs.\ Equiformer) converge to structurally similar
  representations (via CKA), or does the implementation of symmetry matter?
  \item Does the HNN's explicit energy parameterization produce representations
    more aligned with Coulomb physics than symmetry-constrained models, despite
    comparable output accuracy? \item Do any architectures exhibit grokking (a
    delayed generalization transition), and if so, does the emergence of
    physics-aligned representations coincide with this transition?
\end{enumerate}

\subsection{Architecture Suite}

Five architectures will be trained on the identical Coulomb dynamics dataset
under matched hyperparameter conditions:

\begin{center} \renewcommand{\arraystretch}{1.3} \begin{tabular}{llll} \toprule
  \textbf{Architecture}        & \textbf{Tier}             &
  \textbf{Equivariant?} & \textbf{Key Question} \\ \midrule Hamiltonian NN &
  Physics-explicit          & Partial & Does it learn $H$? \\ EGNN &
  Symmetry-constrained      & Yes     & Baseline physics GNN \\ Attention EGNN &
  Symmetry-constrained      & Yes & Does attention learn 1/$r^2$? \\ Equivariant
  Transformer      & Symmetry-constrained      & Yes     & Same symmetry,
different mechanism \\ Non-equivariant GNN          & Unconstrained & No      &
What breaks without symmetry? \\ \bottomrule \end{tabular} \end{center}

\subsection{Experimental Plan}

\textbf{Phase 1: Probing for Physical Quantities}

For each architecture, extract hidden activations at every layer and train
linear probes to predict: (a) pairwise Euclidean distance $\|x_i - x_j\|$, (b)
charge product $q_i q_j \in \{-1, +1\}$, (c) force magnitude $\|F_{ij}\|$, (d)
force direction $\hat{F}_{ij}$, (e) total potential energy $U = \sum_{i<j} q_i
q_j / \|x_i - x_j\|$.

This produces a layer-by-layer ``physics encoding profile'' for each
architecture, revealing what is encoded, at what depth, and whether the tier of
physics prior determines the profile.

\textbf{Phase 2: Causal Verification}

Probing reveals correlation, not causation. For each identified encoding, we
perform causal ablation: project out the probed direction from the residual
stream / hidden state, re-run inference, and measure the MSE increase. We
additionally perform activation patching by corrupting a single particle's
charge (flip $q_i \to -q_i$) and tracing which layers / components carry this
change to the output trajectory. Together, these establish which physical
representations are causally necessary for correct prediction.

\textbf{Phase 3: Representational Geometry}

Compute centered kernel alignment (CKA) between all pairs of architectures at
matched depths to produce a representational similarity matrix. This reveals
whether the two equivariant GNNs (EGNN, Attention EGNN), the equivariant
transformer, and the HNN develop structurally equivalent or divergent internal
representations despite producing similar outputs.

\textbf{Phase 4: Generalization Dynamics}

Train all models across a range of dataset sizes with extended training budgets.
Track probe accuracy for key physical quantities over training steps, alongside
train and validation MSE, to test for grokking-like transitions. The hypothesis:
for physics-explicit architectures, the emergence of correct physics
representations and generalization co-occur; for symmetry-constrained models,
geometry (distance, direction) is learned early but force magnitude may never be
explicitly encoded.

% -------------------------------------------------------
\section{Significance}

This work makes contributions at the intersection of two active research areas:
physics-informed machine learning and mechanistic interpretability.

\textbf{For physics ML:} The field currently evaluates models by output accuracy
and symmetry preservation. This thesis introduces a third axis of evaluation,
\textit{representational alignment with physical law}, and provides the first
systematic comparison of physics-explicit vs.\ symmetry-constrained vs.\
unconstrained architectures on this axis. The results will inform the principled
selection of architecture type for scientific ML tasks where interpretability
and trustworthiness matter.

\textbf{For interpretability:} The overwhelming majority of mechanistic
interpretability work targets transformer language models. Physics simulation
models are an underexplored domain with a critical advantage: the ground-truth
computation is analytically known. This makes it possible to ask precise,
falsifiable questions (``does layer 3 encode Coulomb potential?'') rather than
open-ended ones (``what does this attention head do?''). Insights about
representation formation in physics models may generalize to other structured
scientific domains.

\textbf{The broader question:} Does baking in known symmetries cause a model to
internalize the corresponding physics, or merely constrain its output
transformations while leaving internal computation opaque? This distinction
matters practically: a model that has internalized Coulomb's law should
generalize to novel charge configurations and particle counts; one that has
learned a correlated shortcut may not.

\bibliographystyle{unsrtnat}
\bibliography{../thesis/references}

\end{document}
